\documentclass[12pt, a4paper]{article}
\usepackage[spanish]{babel}
\usepackage{color,graphicx}
\usepackage{graphics}
\usepackage{amsmath,amssymb,amsfonts,amsthm}
\usepackage[T1]{fontenc}
\usepackage{float}
\usepackage{xcolor}
\usepackage{amsthm}
\usepackage{tikz}
\usepackage{enumerate}
\usepackage{picinpar}
\usepackage{newtxtext,fourier}
\makeatother
\usepackage{thmtools}
\usepackage[framemethod=TikZ]{mdframed}
\mdfsetup{skipabove=1em,skipbelow=1em}

\newcommand{\sqed}{\hfill\includegraphics[height=2ex]{sancoya.png}}

\theoremstyle{definition}
\declaretheoremstyle[
    headfont=\bfseries\color{black!70!black}, bodyfont=\normalfont,
    mdframed={
        linewidth=0.7pt,
        rightline=true, topline=true, bottomline=true,
        linecolor=black, backgroundcolor=black!2!white,
    }
]{thmbox}

\declaretheorem[style=thmbox, name=Definición]{definition}
\declaretheorem[style=thmbox, numbered=no, name=Ejemplo]{example}
\declaretheorem[style=thmbox, name=Proposición]{prop}
\declaretheorem[style=thmbox, name=Teorema, numbered=no]{theorem}
\declaretheorem[style=thmbox, name=Lema]{lemma}
\declaretheorem[style=thmbox, numbered=no, name=Corolario]{corollary}
\declaretheorem[style=thmbox, numbered=no, name=Nota]{note}




\usepackage{geometry}
 \geometry{
 a4paper,
 total={170mm,260mm},
 left=25mm,
 top=30mm,
 }
\begin{document}
\setlength{\parindent}{0cm}
\hoffset-0.46cm
\voffset-1.46cm
\noindent
\hbox to \linewidth{
  \begin{minipage}[t]{40pt}
    \vspace*{0pt}
\includegraphics[scale=0.22]{Graphics/logo.eps}
  \end{minipage}
  \hss
  \begin{minipage}[t]{0.8\linewidth}
    \vspace*{15pt}\centering
     Universidad Nacional de Colombia\\
    \small Facultad de Ciencias\\
    Teoría analítica de números\\
    Tarea V
  \end{minipage}
  \hss
}
\phantom{xd}\\
\small
\textsf{Mateo Andrés Manosalva Amaris}
\normalsize
\dotfill

\begin{theorem}
    Sea $\mu(n)$ la función de Möbius y 
$$
M(x) = \sum_{n \leq x} \mu(n).
$$
Entonces
$$
\frac{M(x)}{x} \to 0 \quad \text{cuando } x \to \infty.
$$
\end{theorem}

Antes de elaborar una prueba veamos cómo se comporta la serie de Dirichlet asociada a $\mu(n)$. Sea
\[
g(s):=\zeta(s)-\frac{1}{s-1}.
\]
Tenemos que $g$ es holomorfa alrededor de $s=1$, podemos escribir
$$
\zeta(1+w)=\frac{1}{w}+g(1+w).
$$
Por tanto,
$$
\frac{1}{\zeta(1+w)}
=
\frac{1}{\dfrac{1}{w}+g(1+w)}
=
\frac{w}{1+w g(1+w)}.
$$

Note que
$
w g(1+w)\to 0$ cuando $w\to 0.$
Luego
$1+w g(1+w)\neq 0.$
Así,
$$
\frac{w}{1+w g(1+w)}
$$
es holomorfa alrededor de $w=0$. Además

$$
\lim_{w\to 0}\frac{1}{\zeta(1+w)}=0.
$$

Por lo tanto, $\dfrac{1}{\zeta(s)}$ tiene una singularidad removible en $s=1$ por el teorema de la singularidad removible de Riemann.
Como $\zeta(s)\neq 0$ en $\Re(s)\geq 1$, luego concluimos que $\dfrac{1}{\zeta(s)}$ no tiene polos en la región $\Re(s)\geq 1$. Más aún, observemos su serie de Taylor alrededor de $s=1$, note que 

$$
\zeta(s)=\frac{1}{s-1}+\gamma + a_1(s-1)+a_2(s-1)^2+\cdots ,
$$
donde \(\gamma\) es la constante de Euler--Mascheroni. Entonces
$$
\zeta(s)
=
\frac{1}{s-1}
\left(
1+\gamma(s-1)+a_1(s-1)^2+a_2(s-1)^3+\cdots
\right).
$$
Por lo tanto
$$
\frac{1}{\zeta(s)}
=
(s-1)\dfrac{1}{\left(
1+\gamma(s-1)+a_1(s-1)^2+O((s-1)^3)
\right)}.
$$

Sea
$$
A=\gamma(s-1)+a_1(s-1)^2+O((s-1)^3).
$$

Entonces
$$
(1+A)^{-1}
=
1-\gamma(s-1)-a_1(s-1)^2+\gamma^2(s-1)^2+O((s-1)^3),
$$
obtenemos
\begin{equation}
    \frac{1}{\zeta(s)}
=
(s-1)-\gamma(s-1)^2+(\gamma^2-a_1)(s-1)^3+O((s-1)^4).
\end{equation}
\begin{theorem}
Sea $A$ una función real definida en $[1,\infty)$ tal que, para algún número $\alpha$,
\[
\int_{1}^{\infty} \frac{A(x)-\alpha x}{x^{2}}\,dx
\]
converge.

Supongamos además que se cumple una de las siguientes condiciones:
\begin{enumerate}
  \item[(i)] $A$ es creciente y no negativa; o
  \item[(ii)] existe otra función $B$ tal que $B$ y $B-A$ son crecientes y no negativas, y además, para algún número $\beta$,
  \[
  \int_{1}^{\infty} \frac{B(x)-\beta x}{x^{2}}\,dx
  \]
  converge.
\end{enumerate}

Entonces,
\[
\frac{A(x)}{x} \longrightarrow \alpha
\qquad \text{cuando } x \to \infty.
\]
\end{theorem}

Queremos aplicar el teorema anterior para $A(x)=\sum_{n\leq x} \mu(n)$, esto se hace en virtud de la segunda condición, ya que podemos tomar

$$B(x)=\sum_{n\leq x}1=x,$$

como $|\mu(n)|\leq 1$, se sigue que $A(x)\leq x$, de esta forma $B(x)$ y $(B-A)(x)$ son crecientes y no negativas, tomando $\beta=1$, la integral

\[
  \int_{1}^{\infty} \frac{B(x)-\beta x}{x^{2}}\,dx
  \]

converge. Bastaría ver que, en efecto

\[
\int_{1}^{\infty} \frac{A(x)}{x^{2}}\,dx
\]

converge, ya que en este caso $\alpha=0$ y esto es precisamente, que $A(x)/x\to 0$ cuando $x\to \infty$. Para esto usamos el siguiente teorema. 

\pagebreak

\begin{lemma}
Sea $f$ una función compleja diferenciable en una región que contiene
$$
\{ s : \Re(s) \geq 1 \},
$$
salvo posiblemente en $s = 1$. Supóngase que:

\begin{itemize}
\item[(i)]
$$
f(s) = \int_1^\infty \frac{A(x)}{x^{s+1}}\,dx, \qquad \Re(s) > 1;
$$

\item[(ii)]
$$
f(s) = \frac{\alpha}{s-1} + g(s),
$$
donde $g$ es diferenciable en $s = 1$;

\item[(iii)] existe una constante $M$ tal que
$$
|A(x)| \leq Mx, \qquad x \geq 1.
$$
\end{itemize}

Entonces,
$$
\int_1^\infty \frac{A(x) - \alpha x}{x^2}\,dx
$$
converge y su valor es $g(1)$.
\end{lemma}

\begin{proof}
    
Note que para $\Re(s)>1$

$$g(s)=\sum_{n=1}^{\infty}\dfrac{\mu(n)}{n^s}=1+\int_1^\infty x^{-s}d(A(x))=s\int_1^\infty\dfrac{A(x)}{x^{s+1}}dx,$$

así pues, sea $f(s)=\dfrac{g(s)}{s}=\dfrac{0}{s-1}+\eta(s)$ con $\eta(s)$ analítica en $s=1$ por (1), ya que vimos que $g(s)=\dfrac{1}{\zeta(s)}$ no tiene polos en $s=1$, además, como vimos $|A(x)|\leq x$, de esto se sigue que 

$$
\int_1^\infty \frac{A(x)}{x^2}\,dx
$$
converge a $\eta(1)=\lim_{s\to 1} \dfrac{s}{\zeta(s)}=0$
\end{proof}

Esto completa la prueba. Es tentador escribir que 

$$
\int_1^\infty \frac{A(x)}{x^2}\,dx=\dfrac{1}{\zeta(1)}
$$
 y concluir de manera automáticamente la convergencia de la integral, sin embargo, esto es impreciso, la expresión integral en (i) solo es válida para $\Re(s)>1$, estudiar el comportamiento en $s=1$, la parte Tauberiana de la prueba, permite dar una prueba rigurosa.
\end{document}